% === Begin Label Data ===
% ID: 867
% 难度星级: 
% 标签: 
% 备注: 
% 组卷引用次数: 0
% === End  Label Data ===

\begin{problem}{2026}{G}{新高考I卷}{11}{解析几何}
已知圆 $ C_1:(x+1)^2+y^2=1 $，圆 $ C_2:(x-1)^2+y^2=1 $，圆 $ C_3:x^2+(y-\sqrt{3})^2=1 $，直线 $ l:y=kx+b $ 与 $ C_1,C_2,C_3 $ 均有两个交点. 记 $ l $ 被 $ C_1,C_2,C_3 $ 截得的弦长分别为 $ s_1,s_2,s_3 $，则 (\hspace{1cm})
\begin{choices}
\choice{{ $ k $ 可以取任意实数 }}
\choice{{满足 $ s_1=s_2=s_3 $ 的直线 $ l $ 共有 $ 3 $ 条}}
\choice{{满足 $ s_1+s_2+s_3=3 $ 的直线 $ l $ 多于 $ 3 $ 条}}
\choice{{当 $ b=0 $ 时，$ s_1+s_2+s_3 $ 的最大值为 $ \dfrac{2\sqrt{21}}{3} $ }}
\end{choices}
\end{problem}

\begin{answer}
\end{answer}

\begin{solutions}
\end{solutions}